\newglossaryentry{ancho de banda}{
    type=main,
    name={Ancho de banda},
    description={Frecuencias ocupadas por una señal},
    sort=ancho de banda
}

\newglossaryentry{aprendizaje no supervisado}{
    type=main,
    name={Aprendizaje no supervisado},
    description={Tipo de aprendizaje que utiliza datos no etiquetados},
    sort=aprendizaje no supervisado
}

\newglossaryentry{aprendizaje por refuerzo}{
    type=main,
    name={Aprendizaje por refuerzo},
    description={Tipo de aprendizaje basado en recompensas},
    sort=aprendizaje por refuerzo
}

\newglossaryentry{aprendizaje supervisado}{
    type=main,
    name={Aprendizaje supervisado},
    description={Tipo de aprendizaje que utiliza datos etiquetados},
    sort=aprendizaje supervisado
}

\newglossaryentry{clasificacion}{
    type=main,
    name={Clasificación},
    description={Proceso que asigna una entrada a una clase específica},
    sort=clasificacion
}

\newglossaryentry{clasificacion binaria}{
    type=main,
    name={Clasificación binaria},
    description={Clasificación con solo dos clases posibles},
    sort=clasificacion binaria
}

\newglossaryentry{clutter}{
    type=main,
    name={Clutter},
    description={Conjunto de ecos producidos por objetos que no se mueven},
    sort=clutter
}

\newglossaryentry{compresion de pulsos}{
    type=main,
    name={Compresión de pulsos},
    description={Técnica encargada de mejorar la resolución sin aumentar la potencia de la señal},
    sort=compresion de pulsos
}

\newglossaryentry{correlacion}{
    type=main,
    name={Correlación},
    description={Medida de similitud entre dos señales},
    sort=correlacion
}

\newglossaryentry{desplazamiento doppler}{
    type=main,
    name={Desplazamiento Doppler},
    description={Variación en la señal producida por los objetos en movimiento},
    sort=desplazamiento doppler
}

\newglossaryentry{deteccion}{
    type=main,
    name={Detección},
    description={Proceso que decide si existe o no un objetivo en un eco},
    sort=deteccion
}

\newglossaryentry{direccion angular}{
    type=main,
    name={Dirección angular},
    description={Ángulo que identifica la ubicación del objetivo respecto al radar},
    sort=direccion angular
}

\newglossaryentry{eco de radar}{
    type=main,
    name={Eco de radar},
    description={Señal reflejada por un objeto por el choque de las ondas de radar},
    sort=eco de radar
}

\newglossaryentry{fenomenologia}{
    type=main,
    name={Fenomenología},
    description={Estudio de los fenómenos físicos que alteran la relación entre la señal de radar y el medio},
    sort=fenomenologia
}

\newglossaryentry{filtrado clutter}{
    type=main,
    name={Filtrado clutter},
    description={Proceso encargado de eliminar los elementos estacionarios de la señal},
    sort=filtrado clutter
}

\newglossaryentry{frecuencia}{
    type=main,
    name={Frecuencia},
    description={Número de oscilaciones de una señal por segundo},
    sort=frecuencia
}

\newglossaryentry{frontera de decision}{
    type=main,
    name={Frontera de decisión},
    description={Superficie que separa distintas clases},
    sort=frontera de decision
}

\newglossaryentry{funcion de perdida}{
    type=main,
    name={Función de pérdida},
    description={Medida del error de un modelo de machine learning},
    sort=funcion de perdida
}

\newglossaryentry{funcion sigmoide}{
    type=main,
    name={Función sigmoide},
    description={Función que transforma valores a un rango entre 0 y 1},
    sort=funcion sigmoide
}

\newglossaryentry{interferencia}{
    type=main,
    name={Interferencia},
    description={Señales no deseadas generadas por otras fuentes},
    sort=interferencia
}

\newglossaryentry{machine learning}{
    type=main,
    name={Machine learning},
    description={Técnica que permite aprender patrones a partir de datos},
    sort=machine learning
}

\newglossaryentry{objetivo}{
    type=main,
    name={Objetivo},
    description={Elemento de interés en el proceso de detección de un radar},
    sort=objetivo
}

\newglossaryentry{procesamiento doppler}{
    type=main,
    name={Procesamiento Doppler},
    description={Proceso para obtener la velocidad de los elementos en movimiento en un eco de radar},
    sort=procesamiento doppler
}

\newglossaryentry{rango}{
    type=main,
    name={Rango},
    description={Distancia entre el radar y el objetivo},
    sort=rango
}

\newglossaryentry{regresion lineal}{
    type=main,
    name={Regresión lineal},
    description={Modelo que obtiene un valor continuo a partir de características},
    sort=regresion lineal
}

\newglossaryentry{regresion logica}{
    type=main,
    name={Regresión logística},
    description={Modelo que estima la probabilidad de pertenencia a una clase binaria},
    sort=regresion logica
}

\newglossaryentry{relacion senal ruido}{
    type=main,
    name={Relación señal-ruido},
    description={Comparación entre la señal útil y el ruido del radar},
    sort=relacion senal ruido
}

\newglossaryentry{resolucion}{
    type=main,
    name={Resolución},
    description={Capacidad del radar para distinguir entre dos objetos},
    sort=resolucion
}

\newglossaryentry{ruido}{
    type=main,
    name={Ruido},
    description={Perturbación aleatoria generada al recibir las ondas del radar},
    sort=ruido
}

\newglossaryentry{velocidad radial}{
    type=main,
    name={Velocidad radial},
    description={Componente de la velocidad en dirección de la línea de vista del radar},
    sort=velocidad radial
}

% Manda imprimir el glosario
\printunsrtglossary[type=main,title={\uazGlossaryName}]
